\chapter{Designated Initializers and Aggregate Refinements}
\label{ch:c20-designated-init}

\begin{quote}
\emph{C++20 finally brings designated initializers --- the \texttt{\{.x = 1, .y = 2\}} syntax C programmers have had since C99 --- to C++, along with refinements to what counts as an aggregate. This chapter covers the syntax and its strict rules, range-\texttt{for} init-statements, array-size deduction in \texttt{new}-expressions, and the aggregate model underneath.}
\end{quote}

Designated initializers make struct construction self-documenting, but C++'s version is intentionally stricter than C's: declaration order only, no nesting in one brace pair, no mixing with positional initializers.

\section{Aggregates: The Foundation}
\label{sec:c20-aggregates}

A class is an aggregate when it has no user-declared/inherited constructors, no private/protected non-static data members, no virtual functions, and no virtual/private/protected bases.

\begin{lstlisting}[language=C++,caption={Aggregates --- what qualifies}]
struct Point  { int x; int y; };                               // aggregate
struct Config { int timeout; int retries; bool verbose; };     // aggregate

struct NotAgg { NotAgg(int); int v; };   // user-declared ctor -> NOT aggregate

Point  p{1, 2};
Config c{30, 3, true};
\end{lstlisting}

\section{Designated Initializers}
\label{sec:c20-designated}

\begin{lstlisting}[language=C++,caption={Designated initializers name each field at the call site}]
struct Point { int x; int y; };
Point p{.x = 1, .y = 2};

struct Config { int timeout; int retries; bool verbose; };
Config c{.timeout = 30, .retries = 3, .verbose = true};
Config d{.timeout = 30};          // retries = 0, verbose = false (value-initialized)
\end{lstlisting}

Omitted members are value-initialized.

\section{The Rules: Order, No Mixing, No Repeats}
\label{sec:c20-designated-rules}

\begin{lstlisting}[language=C++,caption={What is and isn't allowed}]
struct S { int a; int b; int c; };

S ok    {.a = 1, .b = 2, .c = 3};   // OK: declaration order
S skip  {.a = 1, .c = 3};           // OK: skip b (value-initialized)
// S bad1  {.b = 2, .a = 1};        // ERROR: out of order
// S bad2  {.a = 1, 2};             // ERROR: mixed designated + positional
// S bad3  {.a = 1, .a = 2};        // ERROR: repeated designator
\end{lstlisting}

Declaration order is mandatory (skipping allowed, reordering not); no mixing positional; each member at most once.

\section{Designated Initializers vs C99}
\label{sec:c20-vs-c99}

\begin{center}
\begin{tabular}{lll}
\hline
\textbf{Capability} & \textbf{C99} & \textbf{C++20} \\
\hline
\texttt{.field = value} & yes & yes \\
out-of-order designators & yes & no \\
mixed designated + positional & yes & no \\
array designators \texttt{[3] = x} & yes & no \\
nested designators \texttt{.a.b = x} & yes & no (nest braces) \\
\hline
\end{tabular}
\end{center}

\begin{lstlisting}[language=C++,caption={Nested aggregates use nested braces}]
struct Inner { int u; int v; };
struct Outer { Inner in; int w; };
Outer o{.in = {.u = 1, .v = 2}, .w = 3};   // C++20
// Outer o{.in.u = 1, .in.v = 2, .w = 3};  // ERROR in C++20 (legal in C99)
\end{lstlisting}

\section{Init-Statements in Range-Based for}
\label{sec:c20-rangefor-init}

\begin{lstlisting}[language=C++,caption={Init-statement in a range-based for}]
#include <vector>
std::vector<int>& getData();

void process() {
    for (auto& data = getData(); auto& x : data) {
        // 'data' scoped to the loop; fixes dangling-temporary iteration
    }
}
\end{lstlisting}

The canonical use is the dangling-range fix: bind a temporary to a named variable whose lifetime spans the loop.

\section{Array Size Deduction in new-Expressions}
\label{sec:c20-new-array-deduction}

\begin{lstlisting}[language=C++,caption={Array size deduced in a new-expression}]
int*  p = new int[]{1, 2, 3};          // size 3 deduced
char* s = new char[]{'a', 'b', '\0'};  // size 3 deduced
// int* q = new int[3]{1, 2, 3};       // previously required explicit count
delete[] p; delete[] s;
\end{lstlisting}

\section{Aggregate Refinements in C++20}
\label{sec:c20-aggregate-refinements}

\begin{lstlisting}[language=C++,caption={The C++20 aggregate tightening and parenthesized init}]
struct A { A() = default; int x; };   // C++20: NOT an aggregate (declared ctor)
// A a{42};                           // ERROR in C++20 (was OK in C++17)

struct Point { int x; int y; };
Point p1{1, 2};                       // brace init
Point p2(1, 2);                       // C++20: parenthesized aggregate init
auto up = std::make_unique<Point>(1, 2);  // now works for aggregates
\end{lstlisting}

A user-declared constructor (even \texttt{= default}/\texttt{= delete}) now makes a type a non-aggregate; parenthesized aggregate init unblocks \texttt{make\_unique}/\texttt{emplace}.

\section{Professional Insights}
\label{sec:c20-designated-insights}

\textbf{Use designated initializers for config and message structs} --- they prevent argument transposition in wide same-typed structs (critical in HFT).

\textbf{The feature requires an aggregate} --- adding any constructor silently disables every \texttt{\{.field=...\}} site.

\textbf{Lean on the init-statement range-for to kill dangling-range bugs.}

\textbf{Treat C99 designated-init habits as non-portable} --- out-of-order, array, and chained designators are rejected by C++20.

\textbf{Prefer deduced array bounds in \texttt{new[]}} --- counts cannot drift out of sync with the data.
